\section{Reward-Influence Implementation Audit}
\label{sec:appendix-influence}

This appendix expands the batch-level audit summarized in
Section~\ref{sec:study}.

We first validated the reward-influence analysis on an unclipped cartpole
batch from the learner indexed by seed 2040.  Reconstructing the implemented gradient from
Equation~\eqref{eq:normalized-update} has error $2.78\times10^{-17}$.  The
clean gradient norm is $0.196789$, neither gradient nor parameter clipping is
active, and the minimum centered-return norm over the reward box is
$7.993864$.  Across 22 release-gate halfspaces, the largest discrepancy
between an SOCP witness and its computed support is $1.07\times10^{-9}$.  The
minimum robust margin was $4.598004$, so no halfspace could be crossed in this
batch.  The successful physical attack instead used the locked bounded
\(\tanh\) target-shaping rule and accumulated influence across changing
batches.  It was not generated by a one-step second-order cone program (SOCP).

The homogenized clipped formulation is then audited on all 36 poisoned V3
batches.  The old premises make 0/36 eligible; the new formulation yields
verified support witnesses on 22/36, including 18 whose reward boxes contain a
normalization singularity and 20 for which the old global gradient-clipping
exclusion fails.  Every actual update is below its computed support, all
witnesses respect the reward budget, final parameter reconstruction error is
zero, and maximum witness/support error is $3.67\times10^{-5}$.  Ten batches
cannot globally exclude coordinatewise parameter clipping (nine activate it
on the observed edit), and four cone-optimal directions have no recovered
witness above the learner's normalization tolerance; they remain outside the
exact implementation audit.  The clipped quadrotor trajectories have not
received this separate parameter-clipping audit and remain empirical.

Exact batch support did not explain the successful multi-batch attack.
The locked \(\tanh\) update advances in the target direction on only 6/20
eligible batches with positive oracle increment, and its preregistered
same-trajectory advantage reaches 22/36 rather than the required 27/36.  The
original exact-support audit emitted 0/12 edits under the earlier singularity
criterion.  A diagnostic rerun on the development realization used the revised
solver and emitted nonzero edits in 9/12 batches.  Its accepted raw snapshot
still caused 0/24 release failures.  The exact-support result therefore serves
as a batch audit.  The physical experiment evaluates the fixed target and
bounded attacker without establishing joint gate optimality of the \(\tanh\)
rule.
